%% Example CV - Magnus Strandgaard
%% Comprehensive overview of all competencies, experience, and achievements.
%% Use as a master reference when tailoring role-specific CVs.
%%
%% Compile with: lualatex main_example.tex
%% (pdflatex often fails on modern MiKTeX with fontawesome5 font-expansion errors.)

\documentclass[11pt,a4paper,sans]{moderncv}
\moderncvstyle{banking}
\moderncvcolor{blue}

% Force both first and last name AND section headings to render in moderncv
% blue (color1). Default banking on lualatex+MiKTeX leaves these black, which
% looks inconsistent with the rest of the blue accent scheme.
\renewcommand*{\firstnamestyle}[1]{{\fontsize{34}{36}\bfseries\upshape\color{color1}#1}}
\renewcommand*{\lastnamestyle}[1]{{\fontsize{34}{36}\bfseries\upshape\color{color1}#1}}
\renewcommand*{\sectionstyle}[1]{{\sectionfont\color{color1}#1}}

\usepackage[utf8]{inputenc}
\usepackage{hyperref}
\hypersetup{
    colorlinks=true,
    linkcolor=blue,
    filecolor=magenta,
    urlcolor=blue,
    pdftitle={Magnus Strandgaard - CV},
    pdfpagemode=FullScreen,
}
\usepackage[scale=0.80]{geometry}
\usepackage{import}

% personal data
\name{Magnus}{Strandgaard}
\address{Copenhagen, Denmark}{}{}
\phone[mobile]{+45 93 99 31 24}
\email{magnus@strandgaard.dev}
\extrainfo{\href{https://linkedin.com/in/magnusstrandgaard}{LinkedIn}, \href{https://github.com/Strandgaard96}{GitHub}}

\begin{document}

\makecvtitle

% ============================================================
%     PROFILE STATEMENT
% ============================================================

\vspace{6pt}
\small{Software Engineer combining a rigorous scientific background with enterprise cloud experience. PhD in computational chemistry with a track record of developing end-to-end Python/ML software workflows for generative molecular design and property prediction, paired with enterprise cloud experience (AWS, Terraform, CI/CD) supporting a 1M+ user application. Driven by a strong passion to build computational systems that accelerate the green energy transition and healthcare.}

% ============================================================
%     CORE COMPETENCIES
% ============================================================

\section{Core Competencies}
\vspace{1pt}
\begin{itemize}

\item \textbf{Python \& Software Engineering}: Python (Expert), FastAPI, REST APIs, Pytest, Git, Docker, CI/CD

\item \textbf{Cloud \& Infrastructure}: AWS (Full Stack), Terraform, Azure DevOps, Jenkins, Linux, Ansible, Proxmox

\item \textbf{Machine Learning}: PyTorch, Scikit-learn, graph neural networks, generative models (genetic algorithms, VAEs), MLOps (WandB, MLflow)

\item \textbf{Cheminformatics \& Computational Chemistry}: RDKit, molecular embeddings, QSPR/QSAR, quantum chemistry (ORCA, xTB, VASP, ADF), HPC/SLURM administration

\item \textbf{AI-Assisted Development}: Claude Code, Gemini CLI, GitHub Copilot - building and shipping production software with AI coding agents

\end{itemize}

% ============================================================
%     PROFESSIONAL EXPERIENCE
% ============================================================

\section{Professional Experience}
\vspace{3pt}
\begin{itemize}

% --- Most Recent Role ---
\item{\cventry{2025--Present}{Linux Operations Engineer}{Netcompany}{Copenhagen, DK}{}{\vspace{1pt}
\begin{itemize}
    \item Maintained and extended AWS infrastructure for a 1M+ user application across 11 environments (10 dev replicas + production), defining compute, networking, and data resources as code in Terraform
    \item Became the team's lead Terraform developer within 3 months of joining; led zero-downtime migration of an enterprise solution (1,000+ live resources) to managed Terraform state
    \item Developed and maintained automated CI/CD pipelines in Azure DevOps and Jenkins with strict policy checks and testing
    \item Engineered a Python data-processing pipeline replacing a legacy service, reducing cloud spend by ~\$3,500/month while ensuring secure data handling
\end{itemize}}}

\vspace{3pt}

% --- Previous Role ---
\item{\cventry{2024--2025}{Postdoctoral Researcher}{Dept. of Chemistry, University of Copenhagen}{Copenhagen, DK}{}{\vspace{1pt}
\begin{itemize}
    \item Designed and operated data pipelines across 4 HPC clusters to ingest, curate, and validate 100K+ molecular records from the Cambridge Structural Database
    \item Built predictive ML workflows combining molecular fingerprints and graph neural networks, with accuracy gains over baseline models; orchestrated distributed GPU training via SLURM and WandB
    \item Developed a novel synthetic accessibility scoring tool in Python with C++ bindings for generative molecular design pipelines
\end{itemize}}}

\vspace{3pt}

% --- Earlier Role ---
\item{\cventry{2021--2024}{PhD Researcher}{Dept. of Chemistry, University of Copenhagen}{Copenhagen, DK}{}{\vspace{1pt}
\begin{itemize}
    \item Developed end-to-end in silico molecular discovery pipelines coupling quantum chemistry with generative ML models (genetic algorithms, variational autoencoders) for catalyst design, screening thousands of candidates with RDKit, Pandas, and SQLite
    \item Collaborated with the University of Oslo theoretical chemistry group to implement a VAE for molecular inverse design, training PyTorch models on HPC GPUs
    \item Served as HPC cluster administrator for the physical chemistry department (50+ active users); created and maintained 4 open-source research repositories; mentored 60+ undergraduate students as teaching assistant
\end{itemize}}}

\vspace{3pt}

% --- Earliest Role ---
\item{\cventry{2021--2021}{Research Assistant}{Dept. of Energy, Technical University of Denmark}{Kongens Lyngby, DK}{}{\vspace{1pt}
\begin{itemize}
    \item Implemented a Message Passing Neural Network (MPNN) in PyTorch to predict forces and energies of magnesium battery cathode candidates, enabling molecular dynamics simulations of next-generation energy-storage materials
\end{itemize}}}

\end{itemize}

% ============================================================
%     EDUCATION
% ============================================================

\section{Education}
\vspace{1pt}
\begin{itemize}

\item{\cventry{2021--2024}{PhD in Computational Chemistry}{University of Copenhagen}{Copenhagen, DK}{}{\vspace{1pt}
Generative ML (genetic algorithms, variational autoencoders) for catalyst design and transition-metal complex inverse design, coupled with quantum chemistry.
}}

\vspace{3pt}

\item{\cventry{2018--2021}{MSc in Physics and Nanotechnology}{Technical University of Denmark}{Kongens Lyngby, DK}{}{\vspace{1pt}
Physics and nanotechnology fundamentals.
}}

\vspace{3pt}

\item{\cventry{2015--2018}{BSc in Physics and Nanotechnology}{Technical University of Denmark}{Kongens Lyngby, DK}{}{\vspace{1pt}
Physics and nanotechnology fundamentals.
}}

\end{itemize}

% ============================================================
%     LANGUAGES
% ============================================================

\section{Languages}
\vspace{1pt}
\begin{itemize}
\item Danish (native), Norwegian (native), English (fluent), Spanish (basic).
\end{itemize}

% ============================================================
%     PUBLICATIONS (optional)
% ============================================================

\section{Publications}
\vspace{1pt}
\begin{itemize}
\item Strandgaard M, Linjordet T, Kneiding H, Burnage AL, Nova A, Jensen JH, Balcells D (2025). A Deep Generative Model for the Inverse Design of Transition Metal Ligands and Complexes. JACS Au. \href{https://doi.org/10.1021/jacsau.5c00242}{DOI link}
\item Strandgaard M, Seumer J, Jensen JH (2024). Discovery of molybdenum-based nitrogen fixation catalysts with genetic algorithms. Chemical Science. \href{https://doi.org/10.1039/D4SC02227K}{DOI link}
\item Strandgaard M, Seumer J, Benediktsson B, Bhowmik A, Vegge T, Jensen JH (2023). Genetic algorithm-based re-optimization of the Schrock catalyst for dinitrogen fixation. PeerJ Physical Chemistry. \href{https://doi.org/10.7717/peerj-pchem.30}{DOI link}
\item Rasmussen MH, Strandgaard M, Seumer J, Hemmingsen LK, Frei A, Balcells D, Jensen JH (2025). SMILES All Around: Structure to SMILES conversion for Transition Metal Complexes. Journal of Cheminformatics. \href{https://doi.org/10.1186/s13321-025-01008-1}{DOI link}
\end{itemize}

% ============================================================
%     HONORS AND AWARDS (optional)
% ============================================================

\section{Honors and Awards}
\vspace{1pt}
\begin{itemize}
\item{No awards on file yet -- add if applicable.}
\end{itemize}

% ============================================================
%     REFERENCES
% ============================================================

\section{References}
\vspace{1pt}
\begin{itemize}
\item{Available upon request.}
\end{itemize}

\end{document}
